% ==============================================================================
% ================================ AUFGABE 6 ===================================
% ==============================================================================

%\chapter{Aufgabe 6: Arduino – Ampelsteuerung}
%\label{chap:aufgabe6}
%\thispagestyle{fancy}

\section{Aufgabe 6: Arduino – Ampelsteuerung}
In diesem Kapitel wird ein Arduino-Sketch entwickelt, der eine einfache Ampelsteuerung mit vier Zuständen realisiert.  
Die Ampel durchläuft zyklisch die Phasen Rot, Rot-Gelb, Grün und Gelb, wobei zwischen jedem Zustand eine Pause von 5 Sekunden eingelegt wird.  
Das Programm läuft in einer Endlosschleife.
Die Umsetzung erfolgt mit der Arduino-IDE gemäß Studienheft MRT13 \cite{3}.
\subsection{Arduino-Sketch}
Der folgende Sketch steuert die Ampel mithilfe von drei LEDs (\textit{Rot}, \textit{Gelb}, \textit{Grün}), die an den digitalen Pins 2, 3 und 4 angeschlossen sind.

\refstepcounter{mylisting}
\begin{minipage}[t]{0.48\textwidth}
    \lstinputlisting[
        language={C},
        style={arduino},
        firstline=1,
        lastline=31,
        frame=tb,
        numbers=left,
        title=\empty,
    ]{asset/code/Aufgabe6.tex}
\end{minipage}
\hfill
\begin{minipage}[t]{0.48\textwidth}
    \lstinputlisting[
        language={C},
        style={arduino},
        firstline=32,
        lastline=59,
        frame=tb,
        numbers=left,
        title=\empty,
    ]{asset/code/Aufgabe6.tex}
\end{minipage}
\begin{minipage}{\linewidth}
    \captionof{lstlisting}{Arduino-Sketch zur Ampelsteuerung (\textit{Quelle: Eigene Darstellung})}
    \label{lst:ampel}
\end{minipage}
\subsection{Schaltungsaufbau}
Die folgende Tabelle zeigt die Pinbelegung der LEDs:
\begin{table}[h]
\centering
\begin{tabular}{c|c|c}
\toprule
\textbf{LED} & \textbf{Arduino-Pin} & \textbf{Farbe} \\
\midrule
Rot   & 2 & Rot \\
Gelb  & 3 & Gelb \\
Grün  & 4 & Grün \\
\bottomrule
\end{tabular}
\caption{Pinbelegung der Ampel-LEDs (\textit{Quelle: Eigene Darstellung})}
\label{tab:ampel_pins}
\end{table}
Jede LED wird über einen Vorwiderstand (\textit{ca. 220~$\Omega$}) mit GND verbunden.

% ==============================================================================
% ================================ AUFGABE 6 ===================================
% ==============================================================================